\documentclass[11pt]{amsart}

\usepackage[a4paper,margin=1in]{geometry}
\usepackage[T1]{fontenc}
\usepackage{lmodern}
\usepackage{microtype}
\usepackage{amsmath,amssymb,amsthm,mathtools}
\usepackage[colorlinks=true,linkcolor=blue,citecolor=blue,urlcolor=blue]{hyperref}

\numberwithin{equation}{section}

\DeclareMathOperator{\tridiag}{tridiag}
\DeclareMathOperator{\diag}{diag}
\DeclareMathOperator{\spec}{spec}
\DeclareMathOperator{\dist}{dist}

\newcommand{\C}{\mathbb C}
\newcommand{\R}{\mathbb R}
\newcommand{\eps}{\varepsilon}
\newcommand{\ii}{\mathrm i}
\newcommand{\A}{A}
\newcommand{\B}{B}
\newcommand{\I}{I}
\newcommand{\J}{J}
\newcommand{\HH}{H}

\newtheorem{theorem}{Theorem}[section]
\newtheorem{mainthm}[theorem]{Main Theorem}
\newtheorem{lemma}[theorem]{Lemma}
\newtheorem{proposition}[theorem]{Proposition}
\newtheorem{corollary}[theorem]{Corollary}
\newtheorem{definition}[theorem]{Definition}
\newtheorem{conjecture}[theorem]{Conjecture}
\newtheorem{remark}[theorem]{Remark}

\title[Connectedness thresholds for Toeplitz pseudospectra]
{Real-Axis Barriers and Connectedness Thresholds for Pseudospectra of a Nonnormal Tridiagonal Toeplitz Family}

\author{Anonymous}

\subjclass[2020]{15A18, 15B05, 47A10, 65F15, 81Q12}
\keywords{Pseudospectrum, nonnormal matrix, tridiagonal Toeplitz matrix, least singular value, connectedness, Hatano--Nelson model, non-Hermitian skin effect}

\begin{document}

\begin{abstract}
We study the Euclidean pseudospectra of the finite nonnormal tridiagonal Toeplitz matrices
\[
A_n(a)=\tridiag(a,0,1),\qquad 0<a<1,\qquad n\ge2.
\]
The spectrum is real because \(A_n(a)\) is diagonally similar to a symmetric Jacobi matrix, but its Euclidean pseudospectra are genuinely nonnormal.  We introduce the real-axis barrier
\[
\tau_n(a)
=
\max_{1\le k\le n-1}
\max_{\lambda_{n,k}(a)\le x\le \lambda_{n,k+1}(a)}
s_{\min}(xI_n-A_n(a)),
\]
where \(\lambda_{n,1}(a)<\cdots<\lambda_{n,n}(a)\) are the eigenvalues of \(A_n(a)\).  The barrier gives the exact connectedness criterion
\[
\sigma_\eps(A_n(a))\ \text{is connected}
\quad\Longleftrightarrow\quad
\eps>\tau_n(a),
\]
provided no off-real bridge occurs before the corresponding real gap closes.  We prove this reduction under the vertical section monotonicity property
\[
s_{\min}((x+\ii y)I_n-A_n(a))
\ge
s_{\min}(xI_n-A_n(a)).
\]
We also give explicit two-sided bounds for \(\tau_n(a)\).  The lower bound is expressed through the optimal moment naturally produced by the Green-kernel estimate,
\[
S_2(a)=\sum_{m=1}^{\infty}m^2a^m=\frac{a(1+a)}{(1-a)^3}.
\]
Consequently,
\[
c_2(a)(n+1)a^{n/2}
\le
\tau_n(a)
\le
(n+1)a^{n/2},
\qquad
c_2(a)=\frac{\sqrt8(1-a)^3}{3\pi^2(1+a)}.
\]
These estimates imply sharp logarithmic bounds for the connectedness threshold.  Finally, we explain the interpretation for the open-boundary Hatano--Nelson chain.
\end{abstract}

\maketitle

\section{Introduction}

For \(M\in\C^{n\times n}\), its \(\eps\)-pseudospectrum is
\[
\sigma_\eps(M)
=
\left\{
z\in\C:\|(zI_n-M)^{-1}\|_2>\eps^{-1}
\right\},
\]
with the convention that the resolvent norm is \(+\infty\) on \(\sigma(M)\).  Equivalently,
\[
\sigma_\eps(M)
=
\left\{
z\in\C:s_{\min}(zI_n-M)<\eps
\right\}.
\]
Pseudospectra are central in nonnormal matrix analysis because eigenvalues alone may fail to describe resolvent growth and spectral instability; see Trefethen \cite{Trefethen1997} and Trefethen--Embree \cite{TrefethenEmbree2005}.

This article concerns
\[
A_n(a)
=
\tridiag(a,0,1)
=
\begin{bmatrix}
0 & 1 \\
a & 0 & \ddots \\
& \ddots & \ddots & \ddots \\
&& \ddots & 0 & 1 \\
&&& a & 0
\end{bmatrix},
\qquad 0<a<1.
\]
Although \(A_n(a)\) is nonnormal, it is diagonally similar to the symmetric matrix
\[
\sqrt a\,\tridiag(1,0,1).
\]
Hence its eigenvalues are real and explicit.  However, pseudospectra in the Euclidean norm are not invariant under nonunitary similarity.  Thus the topology of \(\sigma_\eps(A_n(a))\) reflects nonnormality.

For fixed \(a\in(0,1)\), define the real-axis least singular value
\[
\mu_n(x;a):=s_{\min}(xI_n-A_n(a)),
\qquad x\in\R.
\]
Let
\[
\lambda_{n,1}(a)<\lambda_{n,2}(a)<\cdots<\lambda_{n,n}(a)
\]
be the eigenvalues of \(A_n(a)\).  We define the \(k\)-th gap height by
\[
\beta_{n,k}(a)
:=
\max_{\lambda_{n,k}(a)\le x\le\lambda_{n,k+1}(a)}
\mu_n(x;a),
\qquad k=1,\dots,n-1,
\]
and the largest real-axis barrier by
\[
\tau_n(a):=\max_{1\le k\le n-1}\beta_{n,k}(a).
\]
This quantity is the obstruction to connectedness along the real axis.

The desired finite-dimensional monotonicity statement is
\[
\tau_{n+1}(a)\le\tau_n(a),
\qquad 0<a<1.
\]
If true, it immediately implies that connectedness, once it appears, never disappears as \(n\) increases.  Because this comparison is the hard finite-\(n\) point, we separate unconditional estimates from conditional threshold consequences.

\section{Elementary structure}

Let
\[
B_n(z):=zI_n-A_n(a),
\qquad
s_n(z):=s_{\min}(B_n(z)).
\]
Let \(J_n\) be the reversal matrix,
\[
J_ne_j=e_{n+1-j}.
\]
Then
\[
A_n(a)^TJ_n=J_nA_n(a).
\]
Consequently, for real \(x\),
\[
H_n(x):=J_nB_n(x)=J_n(xI_n-A_n(a))
\]
is real symmetric.  Since \(J_n\) is orthogonal,
\begin{equation}\label{eq:real-axis-smin}
s_n(x)
=
s_{\min}(H_n(x))
=
\min_{\rho\in\sigma(H_n(x))}|\rho|.
\end{equation}

The spectrum of \(A_n(a)\) is explicit.  Put
\[
\Delta_n:=\diag(1,a^{1/2},a,\dots,a^{(n-1)/2}).
\]
Then
\[
\Delta_n^{-1}A_n(a)\Delta_n
=
\sqrt a\,\tridiag(1,0,1).
\]
Therefore
\begin{equation}\label{eq:toeplitz-spectrum}
\sigma(A_n(a))
=
\left\{
2\sqrt a\cos\frac{j\pi}{n+1}:j=1,\dots,n
\right\}.
\end{equation}
We use the increasing order
\[
\lambda_{n,k}(a)
=
2\sqrt a\cos\frac{(n+1-k)\pi}{n+1},
\qquad k=1,\dots,n.
\]
Thus
\[
\lambda_{n,1}(a)<\lambda_{n,2}(a)<\cdots<\lambda_{n,n}(a).
\]

\section{The real-axis barrier}

\begin{definition}
For \(n\ge2\), \(0<a<1\), and \(1\le k\le n-1\), set
\[
G_{n,k}(a)
:=
[\lambda_{n,k}(a),\lambda_{n,k+1}(a)].
\]
Define
\[
\beta_{n,k}(a)
:=
\max_{x\in G_{n,k}(a)}
s_{\min}(xI_n-A_n(a)),
\]
and
\[
\tau_n(a)
:=
\max_{1\le k\le n-1}\beta_{n,k}(a).
\]
\end{definition}

The maximum exists because \(x\mapsto s_{\min}(xI_n-A_n(a))\) is continuous and each \(G_{n,k}(a)\) is compact.  Since \(xI_n-A_n(a)\) is singular exactly at the endpoints of \(G_{n,k}(a)\), one has
\[
\beta_{n,k}(a)>0.
\]

\begin{proposition}[Real-axis connectedness criterion]\label{prop:real-axis-connectedness}
Let
\[
E_{n,\eps}(a)
:=
\{x\in\R:s_{\min}(xI_n-A_n(a))<\eps\}.
\]
Then
\[
E_{n,\eps}(a)\ \text{is connected}
\quad\Longleftrightarrow\quad
\eps>\tau_n(a).
\]
\end{proposition}

\begin{proof}
The function
\[
x\mapsto\mu_n(x;a):=s_{\min}(xI_n-A_n(a))
\]
is continuous and nonnegative.  Its zeros are precisely the real eigenvalues
\[
\lambda_{n,1}(a),\dots,\lambda_{n,n}(a).
\]
Therefore \(E_{n,\eps}(a)\) contains a neighborhood of every eigenvalue.

The only possible disconnections of \(E_{n,\eps}(a)\) between eigenvalues occur inside the closed gaps
\[
G_{n,k}(a)=[\lambda_{n,k}(a),\lambda_{n,k+1}(a)].
\]
The \(k\)-th gap is fully filled by the strict sublevel set \(\mu_n<\eps\) if and only if
\[
\mu_n(x;a)<\eps
\qquad
\text{for every }x\in G_{n,k}(a).
\]
Since
\[
\beta_{n,k}(a)=\max_{x\in G_{n,k}(a)}\mu_n(x;a),
\]
this is equivalent to
\[
\eps>\beta_{n,k}(a).
\]
All real gaps are filled if and only if
\[
\eps>\max_{1\le k\le n-1}\beta_{n,k}(a)=\tau_n(a).
\]
The strict inequality is essential because the pseudospectrum is defined by \(s_{\min}<\eps\).  If \(\eps=\tau_n(a)\), then at least one real gap contains a point where \(\mu_n=\eps\), and that point is not included in \(E_{n,\eps}(a)\).
\end{proof}

\section{Planar connectedness under vertical monotonicity}

The preceding proposition concerns the real trace.  To pass from real-axis connectedness to planar pseudospectral connectedness, one needs a vertical section property.

\begin{definition}[Vertical monotonicity property]\label{def:vertical}
We say that \(A_n(a)\) has the vertical monotonicity property if
\[
s_{\min}((x+\ii y)I_n-A_n(a))
\ge
s_{\min}(xI_n-A_n(a))
\qquad
(x,y\in\R).
\]
\end{definition}

\begin{proposition}[Reduction to real trace]\label{prop:planar-reduction}
Assume that \(A_n(a)\) has the vertical monotonicity property.  Then
\[
\sigma_\eps(A_n(a))\ \text{is connected}
\quad\Longleftrightarrow\quad
E_{n,\eps}(a)\ \text{is connected}.
\]
Consequently,
\[
\sigma_\eps(A_n(a))\ \text{is connected}
\quad\Longleftrightarrow\quad
\eps>\tau_n(a).
\]
Moreover, if \(\sigma_\eps(A_n(a))\) is connected, then it is simply connected.
\end{proposition}

\begin{proof}
For fixed \(x\), define
\[
\Omega_x
:=
\{y\in\R:x+\ii y\in\sigma_\eps(A_n(a))\}.
\]
By vertical monotonicity, if \(y\in\Omega_x\) and \(|y'|\le |y|\), then \(y'\in\Omega_x\).  Hence \(\Omega_x\) is either empty or an open interval symmetric about \(0\).  It is nonempty if and only if \(x\in E_{n,\eps}(a)\).

Thus the vertical projection
\[
\pi:\sigma_\eps(A_n(a))\to E_{n,\eps}(a),
\qquad
\pi(x+\ii y)=x,
\]
has connected fibers, and \(E_{n,\eps}(a)\) is embedded in \(\sigma_\eps(A_n(a))\) as the real trace.  Therefore the connected components of \(\sigma_\eps(A_n(a))\) correspond exactly to the connected components of \(E_{n,\eps}(a)\).  The equivalence with \(\eps>\tau_n(a)\) follows from Proposition \ref{prop:real-axis-connectedness}.

Assume now that \(\sigma_\eps(A_n(a))\) is connected.  Then \(E_{n,\eps}(a)\) is an open interval.  Define
\[
\Phi_t(x+\ii y):=x+\ii(1-t)y,
\qquad 0\le t\le1.
\]
If \(x+\ii y\in\sigma_\eps(A_n(a))\), then vertical monotonicity gives
\[
s_{\min}(\Phi_t(x+\ii y)I_n-A_n(a))
\le
s_{\min}((x+\ii y)I_n-A_n(a))
<
\eps.
\]
Hence \(\Phi_t\) remains inside \(\sigma_\eps(A_n(a))\).  Therefore \(\sigma_\eps(A_n(a))\) strongly deformation retracts onto the real interval \(E_{n,\eps}(a)\), and so it is simply connected.
\end{proof}

\section{The finite monotonicity conjecture and threshold consequences}

\begin{conjecture}[Real-axis barrier monotonicity]\label{conj:barrier}
For every \(0<a<1\) and every \(n\ge2\),
\[
\tau_{n+1}(a)\le\tau_n(a).
\]
\end{conjecture}

\begin{remark}
Conjecture \ref{conj:barrier} is the finite comparison principle needed to turn the exact connectedness criterion into a first-hit equals eventual-hit theorem.  It is not a consequence of ordinary Hermitian eigenvalue interlacing, because \(A_n(a)\) is only nonunitarily similar to a Hermitian matrix and the Euclidean singular values are not preserved by this similarity.
\end{remark}

For fixed \(a\in(0,1)\) and \(\eps>0\), define
\[
N_f(a,\eps)
:=
\min\{n\ge2:\sigma_\eps(A_n(a))\ \text{is connected}\},
\]
and
\[
N_e(a,\eps)
:=
\min\{N\ge2:\sigma_\eps(A_n(a))\ \text{is connected for every }n\ge N\}.
\]

\begin{theorem}[Conditional threshold identity]\label{thm:conditional-threshold}
Assume that the vertical monotonicity property holds for all \(n\), and assume Conjecture \ref{conj:barrier}.  Then
\[
N_f(a,\eps)=N_e(a,\eps).
\]
Writing the common value as
\[
N_c(a,\eps):=N_f(a,\eps)=N_e(a,\eps),
\]
one has
\[
N_c(a,\eps)
=
\min\{n\ge2:\eps>\tau_n(a)\}.
\]
\end{theorem}

\begin{proof}
By Proposition \ref{prop:planar-reduction},
\[
\sigma_\eps(A_n(a))\ \text{is connected}
\quad\Longleftrightarrow\quad
\eps>\tau_n(a).
\]
By Conjecture \ref{conj:barrier}, the sequence \(\tau_n(a)\) is nonincreasing.  The upper bound in Theorem \ref{thm:tau-bounds} below implies
\[
\tau_n(a)\le(n+1)a^{n/2}\to0.
\]
Hence the set
\[
\{n\ge2:\eps>\tau_n(a)\}
\]
is a nonempty tail set.  Its first element is both the first connectedness index and the eventual connectedness index.
\end{proof}

\section{Explicit Green kernel}

The quantitative estimates use the inverse of
\[
B_n(x)=xI_n-A_n(a)
\]
for \(x\) away from the spectrum.  Write
\[
q:=\sqrt a.
\]
For
\[
x=2q\cos\theta,
\qquad
\sin((n+1)\theta)\ne0,
\]
the matrix \(B_n(x)\) is invertible.

\begin{lemma}[Green kernel]\label{lem:green}
Let \(N:=n+1\), \(q:=\sqrt a\), and \(x=2q\cos\theta\), with \(\sin\theta\ne0\) and \(\sin(N\theta)\ne0\).  Then
\[
(B_n(x)^{-1})_{ij}
=
\frac{
q^{i-j-1}\sin(i\theta)\sin((N-j)\theta)
}{
\sin\theta\,\sin(N\theta)
},
\qquad i\le j,
\]
and
\[
(B_n(x)^{-1})_{ij}
=
\frac{
q^{i-j-1}\sin(j\theta)\sin((N-i)\theta)
}{
\sin\theta\,\sin(N\theta)
},
\qquad i>j.
\]
Only absolute values will be used; the harmless global sign depends on the convention for the inverse of the Dirichlet Jacobi matrix.
\end{lemma}

\begin{proof}
Let
\[
\Delta_n=\diag(1,q,q^2,\dots,q^{n-1}).
\]
Then
\[
\Delta_n^{-1}A_n(a)\Delta_n=qT_n,
\]
where
\[
T_n=\tridiag(1,0,1)
\]
is the symmetric Dirichlet Jacobi matrix.  Hence
\[
B_n(x)
=
\Delta_n(xI_n-qT_n)\Delta_n^{-1}.
\]
Thus
\[
B_n(x)^{-1}
=
\Delta_n(xI_n-qT_n)^{-1}\Delta_n^{-1}.
\]
For \(x=2q\cos\theta\),
\[
xI_n-qT_n
=
q(2\cos\theta\,I_n-T_n).
\]
The standard Dirichlet Green kernel for \(2\cos\theta\,I_n-T_n\) is
\[
\left(2\cos\theta\,I_n-T_n\right)^{-1}_{ij}
=
\frac{
\sin(\min\{i,j\}\theta)\sin((N-\max\{i,j\})\theta)
}{
\sin\theta\,\sin(N\theta)
},
\]
up to an overall sign.  Multiplying by \(q^{-1}\) and by the diagonal similarity factor \(q^{i-j}\) gives the displayed formulas.
\end{proof}

\section{Sharp two-sided bounds with \texorpdfstring{\(S_2\)}{S2}}

Define
\[
S_2(a):=\sum_{m=1}^{\infty}m^2a^m.
\]
The closed form is
\[
S_2(a)=\frac{a(1+a)}{(1-a)^3}.
\]
Set
\[
c_2(a)
:=
\frac{\sqrt8\,a}{3\pi^2S_2(a)}
=
\frac{\sqrt8(1-a)^3}{3\pi^2(1+a)}.
\]

\begin{theorem}[Two-sided real-axis barrier bounds]\label{thm:tau-bounds}
For every \(0<a<1\) and every \(n\ge2\),
\[
c_2(a)(n+1)a^{n/2}
\le
\tau_n(a)
\le
(n+1)a^{n/2}.
\]
\end{theorem}

\begin{proof}
We first prove the upper bound.  Let
\[
x=2\sqrt a\cos\theta,
\qquad 0<\theta<\pi.
\]
Define
\[
q^{(n)}(\theta)
=
\left(
\sin\theta,\,
a^{1/2}\sin2\theta,\,
a\sin3\theta,\,
\dots,\,
a^{(n-1)/2}\sin n\theta
\right)^T.
\]
Using
\[
2\cos\theta\sin(j\theta)
=
\sin((j+1)\theta)+\sin((j-1)\theta),
\]
one obtains
\[
B_n(x)q^{(n)}(\theta)
=
a^{n/2}\sin((n+1)\theta)e_n.
\]
Therefore
\[
\mu_n(x;a)
\le
\frac{
a^{n/2}|\sin((n+1)\theta)|
}{
\left(\sum_{j=1}^n a^{j-1}\sin^2(j\theta)\right)^{1/2}
}.
\]
Since
\[
\left(\sum_{j=1}^n a^{j-1}\sin^2(j\theta)\right)^{1/2}
\ge |\sin\theta|
\]
and
\[
|\sin((n+1)\theta)|\le(n+1)|\sin\theta|,
\]
we get
\[
\mu_n(x;a)\le(n+1)a^{n/2}
\]
for every \(x\in(-2\sqrt a,2\sqrt a)\).  All real spectral gaps lie in this interval, hence
\[
\tau_n(a)\le(n+1)a^{n/2}.
\]

We now prove the lower bound.  Let
\[
N:=n+1,
\qquad
\delta:=\frac{3\pi}{2N},
\qquad
\theta:=\pi-\delta.
\]
Then
\[
\frac{(n-1)\pi}{n+1}<\theta<\frac{n\pi}{n+1}.
\]
Consequently,
\[
x_*:=2\sqrt a\cos\theta
\]
lies in the leftmost real spectral gap.  Thus
\[
\tau_n(a)\ge \mu_n(x_*;a).
\]
Since \(x_*\notin\sigma(A_n(a))\),
\[
\mu_n(x_*;a)
=
\frac1{\|B_n(x_*)^{-1}\|_2}
\ge
\frac1{\|B_n(x_*)^{-1}\|_F}.
\]

We estimate the Frobenius norm using Lemma \ref{lem:green}.  For \(i\le j\), set
\[
p:=i,\qquad q_0:=N-j,\qquad m:=p+q_0.
\]
Then \(p,q_0\ge1\), \(m\le N\), and
\[
i-j-1=m-N-1.
\]
Moreover,
\[
|\sin(N\theta)|=1,
\qquad
\sin\theta=\sin\delta,
\]
and, since \(0<\delta\le\pi/2\),
\[
\sin\delta\ge \frac{2\delta}{\pi}.
\]
Also,
\[
|\sin(r\theta)|=|\sin(r\delta)|\le r\delta.
\]
Hence for \(i\le j\),
\[
|(B_n(x_*)^{-1})_{ij}|^2
\le
\frac{\pi^2}{4}\,
p^2q_0^2\delta^2\,
a^{m-N-1}.
\]
For \(i>j\), the same upper bound holds after enlarging the estimate, because the corresponding exponent is \(N-m-1\), and
\[
a^{N-m-1}\le a^{m-N-1}
\qquad(0<a<1,\ m\le N).
\]
Therefore
\[
\|B_n(x_*)^{-1}\|_F^2
\le
\frac{\pi^2}{2}\delta^2a^{-N-1}
\sum_{\substack{p,q_0\ge1\\p+q_0\le N}}
p^2q_0^2a^{p+q_0}.
\]
Extending the sum to all \(p,q_0\ge1\) gives
\[
\|B_n(x_*)^{-1}\|_F^2
\le
\frac{\pi^2}{2}\delta^2a^{-N-1}
\left(\sum_{p=1}^{\infty}p^2a^p\right)^2.
\]
Thus
\[
\|B_n(x_*)^{-1}\|_F^2
\le
\frac{\pi^2}{2}\delta^2a^{-N-1}S_2(a)^2.
\]
Since
\[
\delta=\frac{3\pi}{2N},
\]
we obtain
\[
\|B_n(x_*)^{-1}\|_F
\le
\frac{3\pi^2}{\sqrt8}\,
S_2(a)\,
\frac{a^{-(N+1)/2}}{N}.
\]
Therefore
\[
\mu_n(x_*;a)
\ge
\frac{\sqrt8}{3\pi^2S_2(a)}Na^{(N+1)/2}.
\]
Since \(N=n+1\),
\[
Na^{(N+1)/2}
=
(n+1)a^{(n+2)/2}
=
a(n+1)a^{n/2}.
\]
Hence
\[
\tau_n(a)
\ge
\frac{\sqrt8\,a}{3\pi^2S_2(a)}
(n+1)a^{n/2}
=
c_2(a)(n+1)a^{n/2}.
\]
The proof is complete.
\end{proof}

\section{Explicit threshold bounds}

Assume the vertical monotonicity property and define the connectedness threshold by
\[
N_c(a,\eps):=\min\{n\ge2:\sigma_\eps(A_n(a))\ \text{is connected}\}
\]
whenever this first-hit threshold is considered.  By Proposition \ref{prop:planar-reduction},
\[
N_c(a,\eps)=\min\{n\ge2:\eps>\tau_n(a)\}.
\]
This equality is a first-hit formula.  If Conjecture \ref{conj:barrier} holds, it is also the eventual threshold.

Define
\[
N_+(a,\eps)
:=
\min\{n\ge2:(n+1)a^{n/2}<\eps\},
\]
and
\[
N_-(a,\eps)
:=
1+\max\{n\ge2:c_2(a)(n+1)a^{n/2}\ge\eps\},
\]
with the convention that the maximum over the empty set is \(1\).

\begin{corollary}[Two-sided finite-size threshold bounds]\label{cor:Nc-bounds}
Under the vertical monotonicity property,
\[
N_-(a,\eps)\le N_c(a,\eps)\le N_+(a,\eps).
\]
If Conjecture \ref{conj:barrier} holds, the same bounds apply to the common threshold
\[
N_e(a,\eps)=N_f(a,\eps).
\]
\end{corollary}

\begin{proof}
If
\[
(n+1)a^{n/2}<\eps,
\]
then Theorem \ref{thm:tau-bounds} gives
\[
\tau_n(a)<\eps.
\]
Hence \(\sigma_\eps(A_n(a))\) is connected, and so
\[
N_c(a,\eps)\le n.
\]
Taking the least such \(n\) gives
\[
N_c(a,\eps)\le N_+(a,\eps).
\]

Conversely, if
\[
c_2(a)(n+1)a^{n/2}\ge\eps,
\]
then Theorem \ref{thm:tau-bounds} gives
\[
\tau_n(a)\ge\eps.
\]
Thus \(\sigma_\eps(A_n(a))\) is not connected.  Therefore the connectedness threshold must be larger than every such \(n\), which gives
\[
N_c(a,\eps)\ge N_-(a,\eps).
\]
\end{proof}

Let
\[
\gamma:=\frac{|\log a|}{2}.
\]
Then
\[
a^{n/2}=e^{-\gamma n}.
\]
The equation
\[
C(n+1)e^{-\gamma n}=\eps
\]
with \(C>0\) is solved by the negative branch \(W_{-1}\) of Lambert's \(W\)-function:
\[
n
=
-\frac1\gamma
W_{-1}\!\left(-\frac{\gamma\eps e^{-\gamma}}{C}\right)-1.
\]

\begin{corollary}[Small-\(\eps\) scale]\label{cor:small-eps}
For fixed \(0<a<1\),
\[
N_c(a,\eps)
=
\frac{2}{|\log a|}
\left(
\log\frac1\eps+\log\log\frac1\eps
\right)
+
O_a(1),
\qquad
\eps\downarrow0,
\]
under the vertical monotonicity property.  If Conjecture \ref{conj:barrier} holds, the same asymptotic applies to the common eventual/first-hit threshold.
\end{corollary}

\begin{proof}
Apply Corollary \ref{cor:Nc-bounds} with \(C=1\) and \(C=c_2(a)\).  The standard expansion
\[
W_{-1}(-t)
=
\log t-\log|\log t|+O\!\left(\frac{\log|\log t|}{|\log t|}\right),
\qquad t\downarrow0,
\]
gives
\[
N_c(a,\eps)
=
\frac1\gamma
\left(
\log\frac1\eps+\log\log\frac1\eps
\right)
+
O_a(1).
\]
Since \(\gamma=|\log a|/2\), the result follows.
\end{proof}

\section{Physical interpretation: the open-boundary Hatano--Nelson chain}

The matrix \(A_n(a)\) is the finite open-boundary Hatano--Nelson Hamiltonian with asymmetric nearest-neighbor hopping after the normalization
\[
t_R=1,\qquad t_L=a,\qquad 0<a<1.
\]
Writing
\[
a=e^{-2g},\qquad g>0,
\]
one has
\[
\Delta_n^{-1}A_n(a)\Delta_n
=
e^{-g}\tridiag(1,0,1),
\qquad
\Delta_n=\diag(1,e^{-g},e^{-2g},\dots,e^{-(n-1)g}).
\]
This is the open-boundary imaginary-gauge transformation.  It removes the hopping nonreciprocity at the spectral level but is nonunitary, and therefore does not preserve Euclidean pseudospectra.

The right eigenvectors have the form
\[
r^{(j)}_\ell
=
a^{(\ell-1)/2}\sin\frac{j\ell\pi}{n+1},
\qquad
\ell=1,\dots,n,
\]
whereas the left eigenvectors have the reciprocal profile
\[
\ell^{(j)}_\ell
=
a^{-(\ell-1)/2}\sin\frac{j\ell\pi}{n+1}.
\]
Thus right and left eigenvectors are exponentially biased toward opposite edges.  This is the finite-chain non-Hermitian skin effect.

The connectedness threshold has a direct physical interpretation.  At finite resolution \(\eps\), the open-chain pseudospectrum initially consists of separated spectral islands around individual open-boundary eigenvalues.  The threshold is the length scale at which the largest real-axis barrier is below the resolution level, so the islands merge.  The two-sided estimate
\[
c_2(a)(n+1)a^{n/2}
\le
\tau_n(a)
\le
(n+1)a^{n/2}
\]
shows that the barrier is exponentially small in the chain length, with exponent controlled by the nonreciprocity parameter \(g=|\log a|/2\).  Accordingly,
\[
N_c(e^{-2g},\eps)
=
\frac1g
\left(
\log\frac1\eps+\log\log\frac1\eps
\right)
+
O_g(1).
\]
Stronger nonreciprocity lowers the chain length needed for pseudospectral connectedness.  In the reciprocal limit \(a\uparrow1\), the scale diverges, consistent with the Hermitian limit.

\section{Concluding remarks}

The rigorous unconditional contribution of this paper is the real-axis barrier formulation and the sharp two-sided estimate
\[
\tau_n(a)\asymp_a (n+1)a^{n/2}.
\]
The use of \(S_2(a)\) follows from summing the Green-kernel factors exactly:
\[
\sum_{p,q\ge1}p^2q^2a^{p+q}
=
\left(\sum_{p=1}^{\infty}p^2a^p\right)^2
=
S_2(a)^2.
\]
The finite monotonicity conjecture
\[
\tau_{n+1}(a)\le\tau_n(a)
\]
is the remaining comparison principle needed for the unconditional identity
\[
N_e(a,\eps)=N_f(a,\eps).
\]
A proof of this inequality would upgrade the conditional threshold identity into a complete finite-dimensional topology theorem for this nonnormal Toeplitz family.

\begin{thebibliography}{99}

\bibitem{AmmariBarandunCaoDaviesHiltunen2024}
H.~Ammari, S.~Barandun, J.~Cao, B.~Davies, and E.~O. Hiltunen,
Mathematical foundations of the non-Hermitian skin effect,
\emph{Arch. Ration. Mech. Anal.} \textbf{248} (2024), article no.~33.

\bibitem{BottcherGrudsky2005}
A.~B\"ottcher and S.~M. Grudsky,
\emph{Spectral Properties of Banded Toeplitz Matrices},
SIAM, Philadelphia, 2005.

\bibitem{BottcherEmbreeTrefethen2002}
A.~B\"ottcher, M.~Embree, and L.~N. Trefethen,
Piecewise continuous Toeplitz matrices and operators: slow approach to infinity,
\emph{SIAM J. Matrix Anal. Appl.} \textbf{24} (2002), no.~2, 484--489.

\bibitem{CorlessGonnetHareJeffreyKnuth1996}
R.~M. Corless, G.~H. Gonnet, D.~E.~G. Hare, D.~J. Jeffrey, and D.~E. Knuth,
On the Lambert \(W\) function,
\emph{Adv. Comput. Math.} \textbf{5} (1996), 329--359.

\bibitem{HatanoNelson1996}
N.~Hatano and D.~R. Nelson,
Localization transitions in non-Hermitian quantum mechanics,
\emph{Phys. Rev. Lett.} \textbf{77} (1996), no.~3, 570--573.

\bibitem{HornJohnson2013}
R.~A. Horn and C.~R. Johnson,
\emph{Matrix Analysis},
2nd ed.,
Cambridge University Press, Cambridge, 2013.

\bibitem{KunstEdvardssonBudichBergholtz2018}
F.~K. Kunst, E.~Edvardsson, J.~C. Budich, and E.~J. Bergholtz,
Biorthogonal bulk-boundary correspondence in non-Hermitian systems,
\emph{Phys. Rev. Lett.} \textbf{121} (2018), 026808.

\bibitem{NoschesePasquiniReichel2013}
S.~Noschese, L.~Pasquini, and L.~Reichel,
Tridiagonal Toeplitz matrices: properties and novel applications,
\emph{Numer. Linear Algebra Appl.} \textbf{20} (2013), no.~2, 302--326.

\bibitem{OkumaKawabataShiozakiSato2020}
N.~Okuma, K.~Kawabata, K.~Shiozaki, and M.~Sato,
Topological origin of non-Hermitian skin effects,
\emph{Phys. Rev. Lett.} \textbf{124} (2020), 086801.

\bibitem{ReichelTrefethen1992}
L.~Reichel and L.~N. Trefethen,
Eigenvalues and pseudo-eigenvalues of Toeplitz matrices,
\emph{Linear Algebra Appl.} \textbf{162--164} (1992), 153--185.

\bibitem{Trefethen1997}
L.~N. Trefethen,
Pseudospectra of linear operators,
\emph{SIAM Rev.} \textbf{39} (1997), no.~3, 383--406.

\bibitem{TrefethenEmbree2005}
L.~N. Trefethen and M.~Embree,
\emph{Spectra and Pseudospectra: The Behavior of Nonnormal Matrices and Operators},
Princeton University Press, Princeton, 2005.

\bibitem{YaoWang2018}
S.~Yao and Z.~Wang,
Edge states and topological invariants of non-Hermitian systems,
\emph{Phys. Rev. Lett.} \textbf{121} (2018), 086803.

\end{thebibliography}

\end{document}
